\section{Cap tests for the contact identity}\label{cap:section}

We construct the tests used at a switch.  The construction concerns ordinary
relative invariants of fixed pairs, with disconnected domains permitted.
All descendant lines are the ordinary lines at marked points of maps.
Coefficients used to combine tests are external to the Gromov--Witten
operations.  In particular, the positive operator of
\cref{conv:positive-rule} acts on the auxiliary descendant labels, but does
not act on these coefficients.

\subsection{Geometry, contact spaces, and coefficient conventions}
\label{cap:setup}

The geometric realizations and smoothings of the following pairs are given
in \cref{sew:configurations}.  We specify the data needed here.  Write
\[
 C=\PP_D(\cO\oplus N),\qquad
 D_l=\PP_D(N),\qquad D_r=\PP_D(\cO),
\]
using the convention of lines.  The normal bundles of these sections are
$N^{-1}$ and $N$, respectively.  For a class $\beta$ on $C$, put
$\alpha=\pr_*\beta$ and $k_l=D_l\cdot\beta$, $k_r=D_r\cdot\beta$.
The divisor relation gives
\begin{equation}\label{cap:flux}
 k_r-k_l=\int_\alpha c_1(N).
\end{equation}
The left cap is $P=(C,D_r)$.  In the ordinary configuration the right cap is
$Q=(C,D_l)$.  In the blowup configuration,
\[
 U=\Bl_S M,\quad D=\PP_S(N_{S/M}),\quad N=\cO_D(-1),\qquad
 Q=\bigl(\PP_S(N_{S/M}\oplus\cO),\PP_S(N_{S/M})\bigr).
\]
The right end of the uncut chain in the latter configuration is $V=C$,
relative at $D_l$.

The counting degree $q(\beta)$ is integral and nonnegative on all effective
classes used here.  Ordinarily it comes from a relative polarization ample
on the two original ends and restricts on each neck to a pullback of an
ample class on $D$.  In the blowup configuration it is pulled back from an
ample class on $M$, and on bundle pieces from its restriction to $S$.
We also record any compatible extending divisor degrees required in an
application.  The blowup configuration records the degree against the
nonjoining section $D_r$ of the final $V$, separately on its connected
parts if necessary.  This last degree is not charged to inserted caps.
We use the letter $q$ for the corresponding formal variable as well.

For a smooth, possibly disconnected seam $D$, define the contact superspace
\begin{equation}\label{cap:states}
 \mathcal F_a(D)=
 \bigoplus_{\substack{(m_1,m_2,\ldots)\,;\ m_e\geq0\\
                     \sum_e e m_e=a}}
       \bigotimes_{e\geq1}\Sym^{m_e}\HH(D;\C),
 \qquad
 \mathcal F_{\leq H}(D)=\bigoplus_{0\leq a\leq H}\mathcal F_a(D).
\end{equation}
Here $\Sym$ is graded symmetric: odd factors anticommute.  Only finitely
many $m_e$ can be nonzero.  A contact with a disconnected seam also carries
its component label.  In particular, $\mathcal F_0(D)=\C$ contains the
empty profile.  These spaces are finite dimensional for finite $H$.

Let $G_{\h}$ denote the contact gluing form, with its normalization fixed
by the ordinary degeneration formula.  On ordered roots of orders
$e_1,\ldots,e_\ell$, gluing inserts the categorical Poincare coevaluation
and the factor
\[
        (e_1\cdots e_\ell)\h^\ell.
\]
One then performs the graded symmetrizations and finite permutation
quotients prescribed by that formula.  This describes our normalization
without identifying an ordered profile with an unnormalized monomial.
All permutation denominators and contact factors are nonzero, and
Poincare duality implies that $G_{\h}$ is perfect on every
$\mathcal F_{\leq H}(D)$ over $\C((\h))$.  Profiles with different orders
are orthogonal.  No parity involution is added to the categorical
coevaluation.  The exponent of $\h$ follows from
\[
 g_{\mathrm{total}}-1
    =\sum_v(g_v-1)+\#\{\text{glued pairs of roots}\},
\]
where disconnected arithmetic genus is defined by Euler characteristic.
This formula includes the empty domain.

Fix a contact bound $H$. Our goal is to reproduce $G_{\h}(x,y)$ by ordinary
descendants on two detached caps. Only the incoming state $x$ is initially
bounded: the
$Q$ cap must realize every functional on $\mathcal F_{\leq H}(D)$,
whereas $P$ must prepare a prescribed state on the full outgoing space,
cancelling all unwanted profiles. We first invert the fiber contributions.
In the exceptional configuration, further terms of counting degree zero
strictly lower the outgoing contact; these are inverted next, before
correcting positive counting degrees. The outgoing bound will hold for
each ordinary descendant test before cancellation, so it survives the later
Virasoro action on auxiliary labels. This use of triangular cap calculations is related
to the relative-to-absolute reconstruction of
\cite[Theorem~2 and Sections~2.2--2.4]{MaulikPandharipande2006} and
\cite[Theorem~5.15]{HLR2008}. The contact identity needed here will be
proved below, using the fiber scalar of \cite[Theorem~7.1]{HLR2008}.

We fix a counting bound $d$ throughout each construction.  Auxiliary
markings are encoded by commuting variables for even classes and exterior
variables for odd classes; variables also record descendant indices.
Every coefficient in these variables has only finitely many markings.
A \emph{coefficient row} is the relative functional obtained by extracting
one specified auxiliary monomial from the disconnected cap series, leaving
its contact inputs and its degree and genus series uncontracted.
We work coefficientwise in these variables, modulo $q^{d+1}$, with
Laurent series in $\h$ bounded below.  For finitely many chosen auxiliary
variables the auxiliary completion is
$\C((\h))[[\mathbf t_{\mathrm{even}}]]\ot
\bigwedge(\mathbf t_{\mathrm{odd}})$.  If divisor fugacities are retained,
they are invertible; finite support in their exponents at each step will
be proved below.  The argument never requires a lower bound on the
$\h$-valuation uniform over all auxiliary monomials.

\subsection{The fiber block}

Consider, more generally, the cap
\begin{equation}\label{cap:bundle}
 (B,E)=\bigl(\PP_T(L\oplus\cO),\PP_T(L)\bigr),\qquad \rk L=r,
\end{equation}
with zero section $i:T\hookrightarrow B$ disjoint from $E$.
Allowed ordinary insertions are $\tau_m(i_*\gamma)$ with
$\gamma\in\HH(T)$.  Their first Hodge degrees include the Gysin shift
by $r$.  The ordinary right cap has this form with $L=N$; the left cap
has it with $L=N^{-1}$ after tensoring the projective bundle by $N^{-1}$.
The blowdown cap has $L=N_{S/M}$.

\begin{lemma}[Single-root linearization]\label{cap:linearization}
Restrict \eqref{cap:bundle} to classes with zero projection to $T$.
For any finite bound on contact order, the connected genus-zero
single-root generating series have jointly surjective linearization in
the auxiliary variables.  Consequently one can choose finitely many
variables whose linearization is an isomorphism onto their coefficients
in the contact spaces with one root.
\end{lemma}
\begin{proof}
Choose a homogeneous basis of $\HH(T)$ and use the projective bundle
basis $1,H_f,\ldots,H_f^{r-1}$ on $E$, where $H_f=c_1(\cO_E(1))$.
For root order $e\geq1$ and $0\leq b<r$, use the variable corresponding
to
\begin{equation}\label{cap:linear-row}
       \tau_{r(e-1)+b}(i_*\gamma).
\end{equation}
At a possibly different root order $e'$, pushforward by the relative
evaluation has codimension beyond $\gamma$ equal to
\begin{equation}\label{cap:output-codimension}
                         r(e-e')+b.
\end{equation}
Indeed, writing $t=\dim T$, the relative virtual dimension with one
ordinary and one full-contact relative marking is $t+r-1+re'$.
Subtract the insertion codimension $r+\codim\gamma$ and its descendant
exponent, and compare with $\dim E=t+r-1$.  The projection formula
factors out $\gamma$, so a negative value in
\eqref{cap:output-codimension} forces zero even if the degree of
$\gamma$ itself is large.  Thus orders $e'>e$ do not occur.

At $e'=e$, the output is a nonzero scalar times $H_f^b\gamma$ plus
terms of smaller fiber power with additional base degree.  To compute
the scalar, restrict to a point of $T$.  A connected genus-zero domain
with constant base projection has
$H^1(C,\cO_C)\ot T_tT=0$; horizontal deformations supply the base
factor and no horizontal obstruction.  The calculation is therefore
the relative invariant of $(\PP^r,\PP^{r-1})$.  The exact formula is
\begin{equation}\label{cap:HLR}
 \left\langle\tau_{re-1-j}(\pt)\mid H_f^j\right\rangle^{
               (\PP^r,\PP^{r-1})}_{0,e[\mathrm{line}]}
       =\frac{1}{e^{r-j}((e-1)!)^r},
       \qquad 0\leq j\leq r-1.
\end{equation}
This is \cite[Theorem~7.1]{HLR2008}, for one ordinary marking and one
relative marking of full contact $e$, using the ordinary marked-point
cotangent line. To match the expansion convention used here, apply the
AF-to-Li virtual pushforward of \cite[Theorem~1.1.2]{AMWComparison}.
That comparison takes the relative coarse source without contracting
its components \cite[Section~4.1]{AMWComparison}, so it preserves the
ordinary cotangent line.

For this particular scalar, the jet calculation of
\cite[Lemmas~7.2--7.4]{HLR2008} can be performed on Li's algebraic
relative-map space. Impose evaluation at a point $p\notin\PP^{r-1}$,
and denote its ordinary marking by $x$. The first $e-1$ jets at $x$
form an algebraic section of a bundle with successive quotients
$(L_x^{\otimes k})^{\oplus r}$, $1\leq k<e$. Its top Chern class is
$((e-1)!)^r\psi_x^{r(e-1)}$. The component containing $x$ cannot be
constant: stability would give at least two branches, each forced to
contain the unique relative root, contradicting the genus-zero tree.
At a zero of the jet section, a generic hyperplane through $p$ has
intersection multiplicity at least $e$ at $x$, so this component uses
the full degree $e$. Every remaining rubber branch would have to
carry an outer root by contact conservation. The single-root and
genus-zero conditions therefore leave one full-contact attachment;
any remaining unmarked rubber cylinders are unstable. Thus the zero
locus has unexpanded target and maps $w\mapsto a w^e$. Near it the
point-constrained moduli have the smooth polynomial chart
$[(a_1,\ldots,a_e),\ a_e\ne0\,/\C^*]$, with $a_k\in\C^r$.
The jet equations are regular there, so their localized Euler class
meets the virtual cycle with multiplicity one. The zero locus is a
$\mu_e$-gerbe over $\PP^{r-1}$ and $L_x^{\otimes e}=\cO(1)$ on it.
Its remaining integral is $e^{-(r-j)}$, giving \eqref{cap:HLR}.
This argument only compares the stated one-root fiber invariant.

With $j=r-1-b$, the leading scalar is
$e^{-(b+1)}((e-1)!)^{-r}$.

Order the blocks by root order and, within each block, by fiber power.
The preceding vanishing and nonzero diagonal give a triangular matrix
with invertible diagonal blocks.  Corrections with positive base degree
are triangular in fiber power; alternatively their repeated composition
is nilpotent because the base has bounded cohomological degree.
This proves the assertion in both parities.
\end{proof}

\begin{lemma}[Finite cap spanning]\label{cap:span}
For the fiber classes of \eqref{cap:bundle}, finitely many coefficient
rows in the ordinary zero-section variables span the full dual of
$\mathcal F_{\leq H}(E)$ over $\C((\h))$.  Equivalently, after using
$G_{\h}$ to identify states and their duals, they span all states.
The empty profile and odd labels are included.
\end{lemma}
\begin{proof}
Write $Z_{\mathrm{cl}}(\mathbf t,\h)$ for the disconnected factor
formed by connected components without relative markings.  The
relative product rule in \cref{ev:groups} gives the same factor at
every profile, so it may be divided out before selecting rows.
In fiber degree, a component without roots has degree zero: its
intersection with the relative hyperplane is zero.  At $\mathbf t=0$
its disconnected series is $1+O(\h)$, because unmarked constant maps
in genus zero or one are unstable.  It is therefore invertible over
$\C((\h))$; its inverse is well defined coefficientwise in
$\mathbf t$.  At a fixed auxiliary monomial, only finitely many
coefficients enter multiplication or division by this series.

After this normalization every connected component has at least one
root.  For a profile $\lambda$ with $\ell(\lambda)$ roots, its
lowest possible genus exponent is $-\ell(\lambda)$.  Equality holds
exactly when every connected component has one root and genus zero.
Its coefficient is therefore, up to the nonzero finite labeling
factor, the signed product of the connected genus-zero single-root
series associated with the entries of $\lambda$.

Let these single-root coordinates be $f_1(\mathbf t),\ldots,
f_s(\mathbf t)$ in homogeneous bases.  By
\cref{cap:linearization}, choose $s$ auxiliary variables with invertible
super Jacobian and set the others to zero.  The formal inverse function
theorem applied to $f_i-f_i(0)$ identifies the completed super power
series algebras in these variables.  Substitution by the $f_i$ is
therefore injective on the algebra of polynomials in even variables
and exterior variables in odd degree.  Nonzero constant terms in the
even $f_i$ merely translate the polynomial variables.  It follows that
the finitely many signed monomials corresponding to a basis of
$\mathcal F_{\leq H}(E)$ are linearly independent.

For completeness, finite coefficient detection here involves no
compactness assumption on formal series.  Intersect the kernels of
successively more coefficient functionals on the finite-dimensional
span of these monomials.  If their intersection were nonzero, it would
contain a nonzero series with every coefficient zero.  The decreasing
sequence of subspaces stabilizes, so a finite set of coefficients
already has zero kernel.  Choose a square full-rank submatrix of those
coefficient rows.  Multiply each column of the corresponding full
Gromov--Witten matrix by $\h^{\ell(\lambda)}$.  Its constant matrix
is the just chosen invertible coefficient matrix, and the remaining
entries have positive $\h$-valuation.  It is invertible over
$\C[[\h]]$; undoing the column shifts gives a Laurent inverse.

Finally, each of the finitely many normalized rows used above is a
finite linear combination of unnormalized coefficient rows, with
Laurent coefficients coming from $Z_{\mathrm{cl}}^{-1}$.  Enlarge the
row collection by these finitely many contributing original rows.
Their span still has full rank, so a square subset of original rows
has an invertible Laurent matrix.  Thus all rows can be taken to be
genuine tests by finitely many ordinary markings.  Perfectness of
$G_{\h}$ gives the equivalent assertion for inserted states.
\end{proof}

\subsection{Degree bounds and Laurent completion}

We record the finiteness statement needed to deform the preceding
fiber matrices.  Bounds always refer to effective classes occurring
in the relative splitting formula, not to all integral classes with
the same intersection numbers.

\begin{lemma}[Bounded rows]\label{cap:bounded-rows}
Fix a finite list of auxiliary rows and a counting bound $d$.
\begin{enumerate}[label=\textup{(\roman*)}]
\item For $Q$, a bound on its relative contact total bounds its
      effective curve degrees.
\item For $P$, a fixed nonnegative degree $h$ against $D_l$, together
      with nonnegative outgoing contact at $D_r$, bounds its effective
      curve degrees.  Moreover its outgoing contact satisfies
      $a\leq h+C_0q(\beta)$, for a constant $C_0\geq0$ independent of
      the row.
\item Under these bounds every entry of the finite row matrices is a
      Laurent series in $\h$ bounded below.  At each counting order
      there are finitely many curve classes and finitely many
      exponents of additional divisor fugacities.
\end{enumerate}
\end{lemma}
\begin{proof}
In the ordinary configuration the projected degree on $D$ is bounded
by the ample counting degree.  The section relation
\eqref{cap:flux}, together with either specified section degree,
then fixes the remaining numerical fiber degree.  In the blowdown
cap, the relative hyperplane is relatively ample over $S$; a large
multiple of the ample counting class on $S$ plus this hyperplane is
ample.  Its intersection is bounded by $C d+a$, where $a$ is the
relative contact total.  This proves (i).

In the exceptional left cap, $-c_1(N)$ is relatively ample on
$D\to S$.  For a sufficiently large integer $A$,
\[
        A q|_D-c_1(N)
\]
is ample on $D$.  The same argument, or ampleness of $q|_D$ in the
ordinary case, gives $c_1(N)\cdot\alpha\leq C_0q(\alpha)$.
Since $a=h+c_1(N)\cdot\alpha\geq0$, one also has
$c_1(N)\cdot\alpha\geq-h$.  Thus the projected ample degree is
at most $Ad+h$.  The section degree $h$ then fixes the remaining
fiber coordinate on $P$.  An ample class on $P$ has bounded degree,
proving (ii).

A bound on an integral ample degree bounds the homology classes of
effective curves, not just their numerical images.  One way to see
this is to choose a Kahler metric representing that ample class:
the integral of every fixed smooth real two-form over a holomorphic
curve is bounded in absolute value by a constant times its area.
Apply this to finitely many representatives of a basis of real
cohomology, and use the integral homology lattice and its finite
torsion subgroup.  Hence only finitely many integral classes occur.
In particular all extra divisor exponents have finite support.

For a fixed auxiliary row, let $m$ be its number of ordinary marks.
The number of nonconstant connected components is bounded by their
total degree against the chosen integral ample class.  Constant
connected genus-zero components require at least three ordinary
marks if they have no roots, and degree-zero components cannot have
positive total contact.  Their number is consequently bounded by
$m$.  These are the only connected components contributing a negative
power of $\h$.  Thus each entry has a lower bound on its
$\h$-valuation.  For a specified total exponent, genera are bounded;
unmarked constant components of genus at least two contribute
positive powers, while marked genus-one constants consume ordinary
marks.  Consequently the disconnected exponential and the finite
matrix products are well defined coefficientwise.  This also proves
(iii) for finite lists of rows.
\end{proof}

\subsection{Preparing incoming and outgoing tests}

\begin{lemma}[Incoming tests from $Q$]\label{cap:Q}
Modulo $q^{d+1}$, the $Q$ cap realizes every linear functional on
$\mathcal F_{\leq H}(D)$ by finitely many ordinary coefficient rows
with external coefficients having nonnegative $q$-powers and Laurent
$\h$-coefficients.
\end{lemma}
\begin{proof}
At $q=0$, projection to the base of $Q$ is zero: that base is $D$ in
the ordinary configuration and $S$ in the blowdown configuration,
and the counting class there is ample.  Choose the finite invertible
fiber matrix $M_0$ of \cref{cap:span}.  Apply the same rows at all
counting degrees to obtain
\[
       M(q)=M_0+qM_1+\cdots+q^d M_d
                 \pmod{q^{d+1}}.
\]
Its entries are well defined by \cref{cap:bounded-rows}.  If
$B=M_0^{-1}(M(q)-M_0)$, then
\begin{equation}\label{cap:Qinverse}
 M(q)^{-1}=\left(\sum_{j=0}^{d}(-B)^j\right)M_0^{-1}
                  \pmod{q^{d+1}}.
\end{equation}
This is a finite sum, and uses no negative $q$-powers.
The target classes in a $Q$ row are already bounded by the incoming
contact.  No independent fiber-degree selection on the smooth
segment containing $Q$ is required.
\end{proof}

The left cap requires an additional degree selection.  The degree against
its nonjoining section $D_l$ is supported on that cap and extends to the
smoothed segment.  Select its coefficient $h\geq0$.  This is an external
coefficient selection on that smooth target, not an insertion on a
relative root.  The corresponding section continues on a segment of
type $P|Q$ as well.

\begin{lemma}[Outgoing states from $P$]\label{cap:P}
Modulo $q^{d+1}$, every prescribed state supported in
$\mathcal F_{\leq H}(D)$ can be prepared by finite combinations of
$P$ rows with imposed nonnegative left-section degrees.  The
coefficients have nonnegative $q$-powers.  Every summand of the
preparation, including before cancellation, has outgoing total at
most $H+C_0d$ at total counting order at most $d$.
The same bound holds if ordinary auxiliary insertions are replaced,
removed, or assigned to a classical operation while their imposed
degree conditions are retained.
\end{lemma}
\begin{proof}
First take projected class $\alpha=0$.  Then
\eqref{cap:flux} says $a=h$.  The rank-one case of
\cref{cap:span}, restricted to total contact $h$, supplies an
invertible block of ordinary rows, with no output at other contact
totals.  Normalize these rows by that inverse.

If $q(\alpha)=0$ but $\alpha\ne0$ in the exceptional configuration,
$\alpha$ is vertical over $S$.  The line bundle $-N$ is ample on
every fiber, and hence
\begin{equation}\label{cap:exceptional-down}
            c_1(N)\cdot\alpha<0,\qquad a<h.
\end{equation}
Thus the full $q=0$ matrix at bounded imposed flux is block triangular
in contact, with precisely the invertible fiber blocks on the
diagonal.  Correcting successively downward in contact inverts it:
its strictly contact-decreasing part is nilpotent on the finite
range $0,\ldots,H$.  This step includes empty-profile errors.
For example, over a point with rank-two normal bundle, $D=\PP^1$
and $P$ is the Hirzebruch surface $\mathbb F_1$.  A projected class
of degree $m$ has $a=h-m$, so $0\leq m\leq h$; the term $m=h$
indeed has empty boundary.

We now correct in increasing $q$-order.  Suppose a remaining error
has coefficient $q^j$ and contact $a$.  Use imposed flux $h=a$ and
the already inverted $q=0$ triangular block to remove it.  Any new
error of positive additional counting degree $k$ has contact at
most $a+C_0k$, by \cref{cap:bounded-rows}.  Starting with contact
at most $H$ at order zero proves inductively the sharper bound
\begin{equation}\label{cap:weighted-bound}
            a\leq H+C_0j
            \quad\text{for errors corrected at order }q^j.
\end{equation}
At each order the contact-decreasing corrections terminate by
\eqref{cap:exceptional-down}; only $d+1$ counting orders are
inspected.  There are therefore finitely many required imposed
fluxes and finitely many coefficient rows, and only nonnegative
powers of $q$ have been used.

A row introduced with coefficient divisible by $q^j$ has imposed
flux at most $H+C_0j$.  If its cap class has counting degree $k$
and $j+k\leq d$, its outgoing contact is at most
$H+C_0j+C_0k\leq H+C_0d$.  This argument is an estimate on each
row's curve classes; it does not use cancellation or its insertion
values.  Altering auxiliary labels, including removing a label to
a classical factor, leaves it valid as long as the external degree
conditions are unchanged.
\end{proof}

Extra compatible divisor fugacities cause no enlargement to a field of
arbitrary rational functions in those variables.  In the fiber blocks,
a fixed contact fixes the fiber class, so its fugacity is a fixed
monomial for that column.  Factor out these finitely many monomials
before inversion.  At $q=0$ the exceptional corrections decrease
contact and have finite degree support by \cref{cap:bounded-rows};
their nilpotent inverse is finite.  At positive $q$-order the truncated
inversion and the preceding recursion use only finitely many products.
They therefore introduce only finite Laurent support at the prescribed
bounds.  In particular the final nonjoining-section degree is neither
used nor altered by these inverse matrices.

\begin{proposition}[Cap realization of the contact identity]
\label{cap:identity}
Fix $H,d\geq0$ in either configuration.  Replace a seam by a right
cap $Q$ on its left and a left cap $P$ on its right.  There is a
finite combination of ordinary descendant coefficient rows on
these caps, with external joint coefficients and the left-section
coefficient selections described above, which equals the original
contact gluing form on every incoming state in
$\mathcal F_{\leq H}(D)$ and every outgoing state, modulo
$q^{d+1}$.  This is equality on the full cohomologically decorated
relative state spaces, including the empty profile.

The coefficients have nonnegative $q$-powers and bounded-below
Laurent $\h$-expansions.  At the inspected counting orders all row
summands have outgoing total at most $H+C_0d$, even when their
auxiliary labels undergo the modifications in \cref{cap:P}.
The equality holds after tensoring with any fixed external graded
space and separately for each replica.
\end{proposition}
\begin{proof}
Choose a homogeneous basis $u_i$ of $\mathcal F_{\leq H}(D)$ and
its algebraic dual $\epsilon_i$.  Express the original gluing as
\[
         G_{\h}(x,y)
            =\sum_i \epsilon_i(x)\,G_{\h}(u_i,y),
         \qquad x\in\mathcal F_{\leq H}(D),
\]
understood as a categorical contraction in the displayed order.
If another order is chosen, insert its Koszul sign.
By \cref{cap:Q}, finite $Q$ tests realize each $\epsilon_i$ on
the incoming bounded domain, with the new seam gluing included.
By \cref{cap:P}, finite $P$ tests realize the state $u_i$ on the
entire outgoing space: in particular they cancel every extraneous
outgoing profile through order $d$.  Compose these tests with the
actual gluing tensors.  Perfectness of $G_{\h}$ absorbs all
contact orders, permutation denominators, and genus shifts into
the external coefficients.  Thus the two new gluings have exactly
the weight of the original single gluing, rather than its square.
The stated formula proves the desired identity.

All inversions and products used are finite at the given bounds,
with the completions justified in \cref{cap:bounded-rows}.
The outgoing estimate is the summandwise estimate of
\cref{cap:P}.  Equality before contraction with any outside state
makes tensoring and independent replicas immediate.  The empty
basis vector has been included throughout \cref{cap:span,cap:P};
no nonempty-profile qualification is implicit.
\end{proof}

\subsection{Preparation at several switches}

\begin{lemma}[Uniform preparation]\label{cap:uniform}
For any finite ordered list of switches and a fixed counting bound,
the preceding tests can be chosen successively so that they are
valid for every subset of switch replacements.  All effective
degrees occurring in the resulting coefficient calculations are
bounded, after fixing the final nonjoining-section degree in the
blowup configuration.  These bounds persist when auxiliary
markings are modified at any switches while their external degree
selections are retained.
\end{lemma}
\begin{proof}
At the original left end, an ordinary ample counting class bounds
the initial matching flux.  In the blowup case, if $E_U$ is the
exceptional divisor, choose $A$ such that $Aq-E_U$ is ample.
For an effective class on this end,
\[
             E_U\cdot\beta\leq A q(\beta).
\]
Its matching flux is $E_U\cdot\beta\geq0$, so both flux and ample
degree on $U$ are bounded.  In a neck, the increase in matching
flux is at most $C_0$ times its counting degree, by
\eqref{cap:flux}.  Nonnegative counting degrees on the complete
chain have total at most $d$.

Choose the incoming bound for the first switch accordingly.
If that switch is retained, flux continues under the same neck
estimate.  If it is replaced, \cref{cap:identity} gives a bound on
the new outgoing flux before cancellation, depending only on the
chosen incoming bound and $d$.  Take the larger of these two
bounds, propagate to the next switch, and repeat.  Since there
are finitely many switches, this produces finite bounds for all
subset replacements without a circular choice of rows.

For an exceptional neck with two matching seams, let its section
degrees be $k_l,k_r\geq0$.  The already proved flux bounds and
\eqref{cap:flux} give a lower bound on $c_1(N)\cdot\alpha$.
The ample class $Aq|_D-c_1(N)$ consequently bounds the projected
degree; either section degree fixes the remaining fiber coordinate.
On the final $V$, the same reasoning uses its recorded value
$k_r$ and bounded incoming $k_l$, even if the recorded value is
negative.  On inserted ends it uses the imposed $h$ or the
bounded relative contact, as in \cref{cap:bounded-rows}.  Ordinary
ends are bounded by their ample counting degree.  This exhausts
the pieces and proves the effective-degree assertion.

For fixed rows, their total ordinary arity is finite.  The genus
argument in \cref{cap:bounded-rows} now applies to every piece,
and to their finite products with inverse-matrix coefficients.
Thus extraction at fixed powers of all replica parameters is
legitimate.  Modifying or removing an auxiliary insertion changes
neither a selected degree nor a section-difference relation; it
can only change a finite marking bound.  The same geometric and
Laurent bounds therefore remain valid.
\end{proof}

The ordinary cap markings in this construction are supported away from
the joining divisor.  In the normal-cone models they extend from their
zero-section sources and may be specialized entirely to the chosen
caps.  Two disjoint caps have independent such extensions even if the
resulting classes on a smoothing satisfy linear relations.  One may
choose homogeneous bases on the sources and expand all inverse tests
accordingly.  These observations will let \cref{sew:transport} preintegrate
ordinary auxiliary markings while recording as positions only those
on which an operator actually acts.
